% © 2025 Ing. David Jaroš — CC BY-NC-ND 4.0
%
% This work is licensed under a Creative Commons Attribution-NonCommercial-NoDerivatives 
% 4.0 International License (CC BY-NC-ND 4.0).
%
% License History: Earlier drafts (up to v0.3) were released under CC BY 4.0. 
% From v0.4 onward, all material is released under CC BY-NC-ND 4.0 to protect 
% the integrity of the theoretical work during ongoing academic development.
%
% See LICENSE.md for full license text.

\ifdefined\INCLUDEMODE
\else
  \documentclass[12pt]{article}
  \usepackage[a4paper, margin=2.5cm]{geometry}
  \usepackage{amsmath, amssymb, amsthm}
  \usepackage{hyperref}
  \usepackage{xcolor}
  \usepackage{mdframed}

  \newtheorem{theorem}{Theorem}[section]
  \newtheorem{lemma}[theorem]{Lemma}
  \newtheorem{proposition}[theorem]{Proposition}
  \theoremstyle{definition}
  \newtheorem{definition}[theorem]{Definition}
  \theoremstyle{remark}
  \newtheorem{remark}[theorem]{Remark}

  \newmdenv[backgroundcolor=yellow!20, linecolor=orange, linewidth=1.5pt]{gapbox}
  \newmdenv[backgroundcolor=green!15, linecolor=green!60!black, linewidth=1.5pt]{resultbox}

  \title{\textbf{Step 5: Lamb Shift at $\varphi = \mathrm{const}$}}
  \author{David Jaroš}
  \date{2026-03-06}

  \begin{document}
  \maketitle

  \begin{abstract}
  We assess the reproducibility of the Lamb shift (the $2S_{1/2}$--$2P_{1/2}$ splitting
  in hydrogen, $\Delta E_\mathrm{Lamb}\approx 1057.8\,\mathrm{MHz}$) in UBT with a
  constant scalar background $\varphi = \mathrm{const}$.
  Assuming Steps~1--4 hold, the three dominant contributions to the Lamb shift
  (electron self-energy, vacuum polarisation, and vertex correction) are structurally
  identical to standard QED.
  \textbf{Verdict: SKETCH — structural argument complete; explicit UV-finite combination
  not computed within UBT (open task L1).}
  \end{abstract}
\fi

\section{Step 5: Lamb Shift at $\varphi = \mathrm{const}$}
\label{sec:lamb_shift}

\textbf{Task reference:} UBT\_v29\_task7, step5\_lamb\_shift.
\textbf{Date:} 2026-03-06.
\textbf{Verdict: SKETCH.}

\subsection{Physical Origin of the Lamb Shift}

The Lamb shift in hydrogen arises at order $\alpha^5 m_e$ (in natural units with
$\hbar = c = 1$) from three one-loop QED corrections to the Coulomb potential:
\begin{enumerate}
  \item \textbf{Electron self-energy} $\Sigma(p)$: the electron emits and reabsorbs a
        virtual photon.
  \item \textbf{Vacuum polarisation} $\Pi_{\mu\nu}(k)$: a virtual $e^+e^-$ pair
        screens the Coulomb field.
  \item \textbf{Vertex correction} $\Lambda^\mu(p',p)$: the electron--photon interaction
        is renormalised.
\end{enumerate}

In standard QED the combination gives
\begin{equation}
  \Delta E_{2S_{1/2}-2P_{1/2}}
  = \frac{\alpha^5 m_e}{6\pi}
    \left(\ln\frac{m_e}{\bar{E}_{2S}} - \frac{5}{8} - \frac{1}{5}\right)
  \approx 1057.8\,\mathrm{MHz},
\end{equation}
where $\bar{E}_{2S}$ is the Bethe mean excitation energy.

\subsection{UBT Contributions at $\varphi = \mathrm{const}$}

Given the results of Steps~1--4:

\begin{enumerate}
  \item \textbf{Photon massless} (Step~1): the Coulomb potential is the standard
        $V(r) = -\alpha/r$.  No Yukawa-type screening from photon mass.

  \item \textbf{Electron mass} $m_e = yv$ (Step~2): the electron self-energy
        $\Sigma(p)$ is computed with the same infrared structure as standard QED
        (massive electron, massless photon).

  \item \textbf{Vacuum polarisation} (Step~3, $\delta B = 0$): the photon self-energy
        receives the standard electron-loop contribution; the $\varphi$-loop is absent
        because $q_\varphi = 0$.

  \item \textbf{Vertex correction} (Step~4): the vertex function $\Lambda^\mu$ at
        $\varphi = \mathrm{const}$ is the standard QED expression.
\end{enumerate}

Since all three loop functions $(\Sigma, \Pi, \Lambda)$ are structurally identical to
standard QED in the $\varphi = \mathrm{const}$ background, the Lamb shift calculation
proceeds identically, and the combination yields the standard $1057.8\,\mathrm{MHz}$
at leading order.

\begin{theorem}[Lamb shift at $\varphi = \mathrm{const}$ — structural]
  If Steps~1--4 hold, then the Lamb shift $\Delta E_{2S_{1/2}-2P_{1/2}}$ in hydrogen
  is reproduced by UBT at $\varphi = \mathrm{const}$ to leading order in $\alpha$.
\end{theorem}

\begin{proof}
  (Sketch.)\ The Lamb shift is a UV-finite combination of the three loop functions.
  Since each loop function is individually identical to the standard QED form in the
  $\varphi = v$ background (Steps~1--4), their combination is also identical.  Therefore
  the Bethe-Salpeter computation goes through unchanged, yielding the standard result.
\end{proof}

\subsection{Open Gap}

The argument above is structural (proof by reduction to Steps~1--4).  An explicit
calculation within UBT — computing $\Sigma$, $\Pi$, $\Lambda$ from the UBT path integral
and verifying their finite combination — has not been done.

\begin{gapbox}
\textbf{Gap L1 (open):}
Explicit computation of the UV-finite Lamb-shift combination from the UBT path integral
at $\varphi = \mathrm{const}$ is lacking.  The structural argument (reduction to Steps~1--4)
is convincing but does not constitute a rigorous derivation within UBT.  Resolving Gap~L1
requires:
\begin{enumerate}
  \item Computing $\Sigma(p)$ from the UBT propagator in the $\varphi = v$ background;
  \item Verifying the UV/IR cancellations in the full Lamb-shift combination;
  \item Comparing the Bethe logarithm against the standard value.
\end{enumerate}
This is an important but not blocking gap: if any of Steps~1--4 were to fail, the Lamb
shift would provide an immediate test.
\end{gapbox}

\subsection{$\varphi$-Corrections to the Lamb Shift}

New contributions at $\varphi \neq 0$ arise from diagrams involving internal $\varphi$
propagators.  These are suppressed relative to the standard Lamb shift by
\begin{equation}
  \frac{\delta(\Delta E)_\varphi}{\Delta E_\mathrm{Lamb}}
  \sim \frac{y_e^2}{16\pi^2} \approx 10^{-14},
\end{equation}
using $y_e \approx 3\times10^{-6}$ (SM electron Yukawa).  This is below current experimental
sensitivity ($\delta(\Delta E)_\mathrm{exp}/\Delta E \sim 10^{-10}$).

\subsection{Summary}

\begin{center}
\begin{tabular}{ll}
\hline
Claim & Lamb shift $1057.8\,\mathrm{MHz}$ reproduced at $\varphi = \mathrm{const}$ \\
Status & \textbf{SKETCH} \\
Structural argument & Reduction to Steps~1--4 (complete) \\
Gap L1 & Explicit UBT path-integral computation lacking \\
$\varphi$-correction & $\sim 10^{-14}$ relative (unobservable) \\
Layer & [L1] \\
\hline
\end{tabular}
\end{center}

\begin{resultbox}
\textbf{Conclusion (Step 5, SKETCH):}
The Lamb shift is structurally reproduced in UBT at $\varphi = \mathrm{const}$, following
directly from Steps~1--4.  The explicit UBT path-integral computation is an open task
(Gap~L1) but is not expected to yield surprises given the structural agreement.
$\varphi$-corrections are of order $10^{-14}$ relative to the Lamb shift and are
unobservable at current precision.
\end{resultbox}

\ifdefined\INCLUDEMODE\else
  \end{document}
\fi
